\section{The transferable mechanism in the Lean source}

\subsection{The arithmetic remainder}

For a positive modulus $q$ and an auxiliary prime $d$, the source uses
\[
 R_{d,q}(x)=
 \sum_{\substack{n\le x\\n\equiv1\pmod q}}
 \Lambda(n)(1-\chi_{-d}(n))
       \left(\frac1n-\frac1x\right)
 -\frac{\log x}{\varphi(q)}-c_{d,q}.
\]
Here $\Lambda(p^k)=\log p$ and is zero otherwise. The constant
$c_{d,q}$ centers the expression at infinity. The field and congruence
conditions guarantee that the quadratic character used in the code
agrees with this notation. For example, the final construction takes
$18480\mid q$ and $d\equiv q/2-1\pmod q$.

The smoothing factor $(1/n-1/x)$ is essential. It inserts
$1/[s(s-1)]$ into the Mellin transform, so zero multiplicities appear
with a controlled sign at real points. The source definitions and
transform identities occur in \path{New/Reduction/Zeros/PrimeError.lean};
the place sums that produce this statistic are in
\path{New/Reduction/PrimeSums.lean} \cite{repo}.

\subsection{Three inputs for a common lattice}

The proof constructs spaces of automorphic forms and compares the same
logarithmic squared covolume $V$, corrected by twice an arithmetic index
$\delta$. Its three key inputs are:

\begin{enumerate}
\item \textbf{Short forms.} Theta derivatives yield integral source
forms whose images have controlled size and divisibility.
\item \textbf{Enough forms.} A second-moment argument shows that enough
independent images are present to occupy a large fraction of the target.
\item \textbf{Integral forms are long.} Heights of determinants of
chains of theta lifts, together with the product formula, provide a
lower bound for the same covolume/index expression.
\end{enumerate}

Write $D$ for target dimension, $n$ for the number of independent images,
$H$ for a shared main term, $m$ for chain length, and $\mathrm{rem}$
for the unnormalized prime remainder. Schematically,
\[
 DH-\text{lower loss}
 \ \le\ V+2\delta\
 \le n(H-m\,\mathrm{rem}+\text{upper loss})
       +\text{index loss}.
\]
The main terms cancel. The exact source result is
\[
 m\,\mathrm{rem}\le f\,a(\log b)^5.
\]
After converting the place sum to $R_{d,q}$, a further factor of
$\log b$ appears. With $m\ge\sqrt x$, $\log b\le2x^\varepsilon$,
and errors controlled at a much larger cutoff $x_1$, the source proves
\[
 R_{d,q}(x)\le(64\varphi(q)+3)x^{-1/2+6\varepsilon}
\]
for large integers $x$. Interpolation and adjustment of $\varepsilon$
give the final all-real-scale one-sided bound
\cite{mainineq,onesided}.

This square-root amplification is a specific arithmetic accomplishment.
Introducing a symbol $m=\sqrt x$ in an unrelated problem would not
reproduce it: the shortness, integrality, spanning, denominator, and
determinant-height estimates must all hold at that chain length.

\subsection{Two distinct uses of twisting}

Twists serve two separate purposes. Inside the arithmetic argument they
control local denominators and compare determinant heights; a good
twist may be selected separately at each coefficient prime. In the
analytic endgame an auxiliary quadratic twist is selected to avoid
zeros in a prescribed rectangle.

The latter is necessary because the transformed remainder detects a
\emph{difference} of zero-multiplicity sums. A zero of an untwisted
function could be canceled by a twisted zero at the same point.
The source's zero-density input selects a $d$ for which the twisted
terms have no zeros in the relevant box. The residual numerator is
then a sum of nonnegative integers. Landau's positivity argument
converts the one-sided upper bound into pole exclusion
\cite{criterion}.

For a full Dedekind-zeta statistic, zero multiplicities can instead
already have the same sign. That removes the auxiliary analytic
zero-avoidance step. It does not remove any twists still needed by
the arithmetic construction.

\subsection{What can be transported and what must be rebuilt}

The scalar determinant comparison, the product-formula logic, and the
Landau transform argument admit clean general statements. In contrast,
the construction of the prime weights and the sufficient supply of
forms are specific to the chosen CM field, level, local types, and
theta chain.

For a new $L$-function, the central design question is therefore:
\emph{which arithmetic object produces its actual prime or prime-ideal
coefficients, with the correct normalization and controlled losses?}
An appealing analogy to theta lifting is insufficient without that
coefficient identity.

The next section changes the analytic interface. It shows that a
successful new arithmetic construction need not necessarily give an
estimate at every individual scale. Suitable averages can supply all
the analytic information needed.
